\begin{tikzpicture}[x=10.000000cm,y=4.363636cm,font=\scriptsize]\draw[otRule!70,line width=.25pt] (0,0)--(0,1.1);\draw[otRule!70,line width=.25pt] (1,0)--(1,1.1);\draw[otRule!70,line width=.25pt] (0,0)--(1.1,0);\draw[otRule!70,line width=.25pt] (0,1)--(1.1,1);\draw[otMuted,->] (0,0)--(1.25,0) node[right] {$x$};\draw[otMuted,->] (0,0)--(0,1.25) node[above] {$y$};\node[below,fill=white,inner sep=.5pt,font=\fontsize{6.5}{7}\selectfont,text=otMuted] at (1,0) {1};\node[left,fill=white,inner sep=.5pt,font=\fontsize{6.5}{7}\selectfont,text=otMuted] at (0,1) {1};\begin{scope}\clip (0,0) rectangle (1.1,1.1);\draw[otNavy,line width=.9pt] plot[domain=0:1,samples=91] (\x,{\x*\x});\geoSolution{\draw[geoReponse](0,0)--(.5,.25)--(1,1);}\end{scope}\end{tikzpicture}
\par\legende{La courbe $y=x^2$ et la ligne polygonale à deux segments. Les graduations donnent l’échelle de chaque axe.}
